% vim: spl=pt
\begin{exercício}{Espectro de um operador unitário}{ex14}
    Mostre que, se \(U\) é um operador unitário, então
    \begin{enumerate}[label=(\alph*)]
        \item \(U\) é uma isometria sobrejetiva (bijetiva); e
        \item \(\sigma(U) \subset \setc{z \in \mathbb{C}}{\abs{z} = 1}\).
    \end{enumerate}
\end{exercício}
\begin{proof}[Resolução de (a)]
  Como \(U\) é unitário, temos \(U \in \bounded(\hilbert)\) e \(U^* = U^{-1}.\) Assim, \(U\) é bijetivo e, em particular, sobrejetivo. Notemos que \(U\) preserva o produto interno, já que
  \begin{equation*}
    \forall x,y \in \hilbert: \inner{Ux}{Uy} = \inner{x}{U^*Uy} = \inner{x}{y}.
  \end{equation*}
  Assim, temos
  \begin{equation*}
    \forall x \in \hilbert: \norm{Ux} = \sqrt{\inner{Ux}{Ux}} = \sqrt{\inner{x}{x}} = \norm{x},
  \end{equation*}
  logo \(U\) é uma isometria.
\end{proof}
\begin{proof}[Resolução de (b)]
  Como \(U\) é unitário, é uma isometria bijetiva, portanto \(\norm{U} = 1\). Ainda, \(U\) é normal já que \(U^* U = U U^*,\) portanto \(r(U) = \norm{U} = 1\) e concluímos que \(\sigma(U) \subset D,\) onde \(D \subset \mathbb{C}\) é o disco fechado unitário. Temos o mesmo resultado para \(U^* = U^{-1}\), isto é, \(\sigma(U^{-1}) \subset D\). Ainda, temos
  \begin{align*}
    \lambda \in \sigma(U) &\iff \lambda \unity - U \notin \mathrm{Inv}(\hilbert)\\
                          &\iff \lambda^{-1} \unity - U^{-1} = -\lambda^{-1} U^{-1} (\lambda \unity - U) \notin \mathrm{Inv}(\hilbert)\\
                          &\iff \lambda^{-1} \in \sigma(U^{-1}).
  \end{align*}
  Disso segue que
  \begin{align*}
    \lambda \in \sigma(U) &\iff \lambda \in \sigma(U) \land \lambda^{-1} \in \sigma(U^{-1})\\
    &\implies\abs{\lambda} \leq 1 \land \abs{\lambda}^{-1} \leq 1\\
    &\iff \abs{\lambda} \leq 1 \land \abs{\lambda} \geq 1\\
    &\iff \abs{\lambda} = 1,
  \end{align*}
  portanto \(\sigma(U) \subset \setc{z \in \mathbb{C}}{\abs{z} = 1}.\)

  Vamos justificar a afirmação de que para um operador normal \(T \in \bounded(\hilbert),\) então \(r(T) = \norm{T}\). Primeiro, consideramos \(S \in \bounded(\hilbert)\) auto-adjunto, e mostramos que \(\norm{S^{(2^n)}} = \norm{S}^{(2^n)}\) para todo \(n \in \mathbb{N}\). Pela propriedade C*, o caso \(n = 1\) é válido, então assumindo válido para \(k \in \mathbb{N}\) temos
  \begin{equation*}
    \norm{S^{(2^{k+1})}} = \norm{S^{(2^k)} S^{(2^k)}} = \norm{[S^{(2^k)}]^* S^{(2^k)}} = \norm{S^{(2^k)}}^2 = [\norm{S}^{(2^k)}]^2 = \norm{S}^{(2^{k+1})},
  \end{equation*}
  portanto válido para \(k + 1.\) Pelo princípio da indução finita, concluímos que a igualdade é válida para qualquer \(n \in \mathbb{N}.\)
  Voltando nossa atenção para o operador normal \(T\), mostramos por indução novamente que \((T^*T)^{(2^n)} = \left[T^{(2^n)}\right]^* T^{(2^n)}\) para todo \(n \in \mathbb{N}.\) De fato, temos
  \begin{equation*}
    (T^*T)^2 = T^*T T^*T = (T^*)^2 T^2 = (T^2)^*T^2
  \end{equation*}
  e assumindo verdadeira para algum \(k \in \mathbb{N}\) temos
  \begin{align*}
    (T^*T)^{(2^{k+1})} &= [(T^*T)^{(2^k)}]^2\\
                       &= \{[T^{(2^k)}]^*T^{(2^k)}\}^2\\
                       &= [T^{(2^k)}]^*T^{(2^k)}[T^{(2^k)}]^*T^{(2^k)}\\
                       &= [T^{(2^k)}]^*[T^{(2^k)}]^*T^{(2^k)}T^{(2^k)}\\
                       &= [T^{(2^{k+1})}]^*T^{(2^{k+1})}
  \end{align*}
  e a validade para todo \(n \in \mathbb{N}\) segue pelo princípio da indução finita. Pela propriedade C*, temos então que
  \begin{equation*}
    \norm{T^{(2^k)}}^2 = \norm{[T^{(2^k)}]^*T^{(2^k)}} = \norm{(T^*T)^{(2^n)}} = \norm{T^*T}^{(2^n)} = \norm{T}^{(2^{n+1})},
  \end{equation*}
  isto é, \(\norm{T^{(2^k)}} = \norm{T}^{(2^k)}.\)
  Consideremos agora a sequência \(\sequence{\norm{T^n}^{\frac1n}}{n\in \mathbb{N}}\) convergente para \(r(T)\). Assim, a subsequência \(\sequence{\norm{T^{(2^k)}}^{\frac{1}{2^k}}}{k \in \mathbb{N}}\) converge para \(r(T)\) e temos
  \begin{equation*}
    r(T) = \lim_{k \to \infty}{\norm{T^{(2^k)}}^{\frac1{2^k}}} = \lim_{k \to \infty}{\norm{T}} = \norm{T},
  \end{equation*}
  isto é, a norma de um operador normal é igual ao seu raio espectral.
\end{proof}
